\chapter{Cambios, limitaciones y cierre}

\section{Resumen de la entrega}

La preparación de cierre incorporó un portal README, documentación por audiencia, guía de usuario,
arquitectura, despliegue/recuperación, plan de pruebas, seguridad, contribución, CODEOWNERS,
plantillas y este manual versionado.

También se alineó la documentación con el comportamiento actual: reportes en PDF/Excel, reapertura
de tickets cerrados, autenticación con access token en memoria y refresh cookie, topología
Vercel/Render, demo protegida y migraciones como paso controlado.

\section{Formatos y compatibilidad}

\begin{itemize}
  \item Reportes: PDF y Excel. CSV no forma parte del alcance actual.
  \item Backend: Python 3.12 en CI, Django 6 y ASGI mediante Daphne.
  \item Frontend: Node 24 en CI, React 19, TypeScript 6 y Vite 8.
  \item Base: PostgreSQL 15 mediante Supabase.
  \item Tiempo real: Channels; Redis requerido para más de un proceso.
\end{itemize}

\section{Limitaciones conocidas}

\begin{enumerate}
  \item no hay contrato OpenAPI publicado;
  \item existe cobertura crítica reproducible, pero no un umbral bloqueante en CI;
  \item no hay un E2E completo automatizado en navegador;
  \item backup y restauración necesitan evidencia periódica con RTO/RPO;
  \item disponibilidad, correo y WebSocket dependen de proveedores externos;
  \item ownership, licencia/cesión y soporte requieren acuerdo escrito;
  \item el commit final necesita PR, CI y tag antes de considerarse release contractual.
\end{enumerate}

\section{Checklist de transferencia}

\begin{tabularx}{\textwidth}{@{}L{0.08\textwidth}Y@{}}
\toprule
\textbf{Estado} & \textbf{Control}\\
\midrule
$\square$ & Commit final publicado, PR aprobado y CI verde.\\
$\square$ & Tag/release y checksums registrados.\\
$\square$ & Smoke por roles, correo y WebSocket documentados.\\
$\square$ & Restore de backup probado y RTO/RPO registrados.\\
$\square$ & Propietario y segundo administrador transferidos para cada plataforma.\\
$\square$ & Secretos entregados por canal seguro y rotados cuando corresponde.\\
$\square$ & Licencia/cesión, soporte y aceptación firmados.\\
$\square$ & Manual PDF archivado junto al tag correspondiente.\\
\bottomrule
\end{tabularx}

\section{Fuentes canónicas}

El mantenimiento debe comenzar en `README.md`, `docs/README.md`, `docs/ARCHITECTURE.md`,
`docs/USER\_GUIDE.md`, `docs/DEPLOYMENT.md`, `docs/TESTING.md`, `SECURITY.md`,
`CONTRIBUTING.md` y `CHANGELOG.md`. Este manual consolida esas fuentes para el cliente; cuando
exista una discrepancia, corrija primero el documento dueño del tema y regenere el PDF.

\ManualNote{No entregue una copia comprimida del workspace. Use un clon limpio, `git archive` o una
release para excluir `.env`, backups, cachés, herramientas locales y material legado.}
